\documentclass[11pt, a4paper]{article}
\usepackage[margin=2.5cm]{geometry}
\usepackage{amsmath, amssymb, amsthm}
\usepackage{bm}
\usepackage{microtype}
\usepackage{hyperref}
\usepackage{xcolor}
\usepackage{booktabs}
\usepackage{tcolorbox}
\usepackage{listings}
\usepackage[T1]{fontenc}
\usepackage{lmodern}

\definecolor{copper}{HTML}{8f3e22}
\definecolor{ivory}{HTML}{f7f3ea}
\hypersetup{colorlinks=true, linkcolor=copper, urlcolor=copper, citecolor=copper}

\theoremstyle{definition}
\newtheorem{definition}{Definition}[section]
\newtheorem{proposition}{Proposition}[section]
\newtheorem{remark}{Remark}[section]

\lstset{language=Python, basicstyle=\small\ttfamily, keywordstyle=\color{copper}\bfseries,
  frame=single, backgroundcolor=\color{ivory}, breaklines=true, columns=fullflexible}

\newcommand{\bx}{\bm{x}}
\newcommand{\bepsilon}{\bm{\epsilon}}
\newcommand{\E}{\mathbb{E}}
\newcommand{\N}{\mathcal{N}}

\title{\color{copper}\textbf{Flow Matching}\\[4pt]
  \large TLH-2 · Thursday Learning Hours · Applied Models Track}
\author{Thursday Learning Hours}
\date{April 17, 2026}

\begin{document}
\maketitle

\begin{tcolorbox}[colback=ivory, colframe=copper, title=\textbf{Core Thesis}]
Diffusion models choose a curved Gaussian path from data to noise. Flow matching chooses
any path — and the simplest choice is a straight line. Train a velocity field network to
predict $\bepsilon - \bx_0$. Integrate the ODE to generate. No SDE, no noise schedule,
no ELBO. Fewer steps. Same quality.
\end{tcolorbox}

\tableofcontents
\newpage

\section{Probability Paths}

Both diffusion and flow matching construct a time-indexed family of distributions
$\{p_t\}_{t \in [0,1]}$ interpolating between $p_0 = p_\text{data}$ and $p_1 = \mathcal{N}(\bm{0}, \bm{I})$.

\textbf{Key insight:} any smooth path can be described by an ODE:
\begin{equation}
  \frac{d\bx}{dt} = v_t(\bx), \quad \bx_0 \sim p_0
  \label{eq:ode}
\end{equation}

The density along this path satisfies the continuity equation:
$\partial_t p_t + \nabla \cdot (v_t p_t) = 0$.

\section{Rectified Flow}

Liu et al.~\cite{liu2022rf} choose the simplest path — linear interpolation:
\begin{equation}
  \bx_t = (1-t)\,\bx_0 + t\,\bepsilon, \quad \bepsilon \sim \mathcal{N}(\bm{0}, \bm{I}), \quad t \in [0,1]
  \label{eq:rf_path}
\end{equation}

The velocity is constant along this path:
\begin{equation}
  v_t(\bx_t) = \frac{d\bx_t}{dt} = \bepsilon - \bx_0
\end{equation}

\subsection{Training Objective}

\begin{equation}
  \boxed{\mathcal{L}_\text{RF} = \mathbb{E}_{t, \bx_0, \bepsilon}\!\left[\|(\bepsilon - \bx_0) - v_\theta(\bx_t, t)\|^2\right]}
\end{equation}

\subsection{Euler Sampling}

Integrate the learned ODE from $t=1$ to $t=0$:
\begin{equation}
  \bx_{t - \Delta t} = \bx_t - \Delta t \cdot v_\theta(\bx_t, t)
\end{equation}

With $n$ steps: $\Delta t = 1/n$. Straighter paths → larger $\Delta t$ → fewer steps → same quality.

\section{Conditional Flow Matching}

Lipman et al.~\cite{lipman2022fm} observe that marginal paths from random $(x_0, \epsilon)$ pairs are
slightly curved (pairs ``cross'' in data space). Conditioning on $x_0$ removes this:

\begin{equation}
  \mathcal{L}_\text{CFM} = \mathbb{E}_{t, \bx_0, \bx_t \sim p_t(\cdot|\bx_0)}\!\left[\|u_t(\bx_t | \bx_0) - v_\theta(\bx_t, t)\|^2\right]
  \label{eq:cfm}
\end{equation}

where the conditional velocity $u_t$ is known in closed form for linear paths.

\subsection{Optimal Transport Coupling}

Pairing each $\bx_0$ with its nearest $\bepsilon$ (minibatch OT) gives even straighter paths,
approaching geodesics in the Wasserstein metric.

\section{Comparison with DDPM}

\begin{center}
\begin{tabular}{lll}
\toprule
Property & DDPM & Rectified Flow \\
\midrule
Path & $\sqrt{\bar\alpha_t}\bx_0 + \sqrt{1-\bar\alpha_t}\bepsilon$ (curved) & $(1-t)\bx_0 + t\bepsilon$ (linear) \\
Target & $\bepsilon$ & $\bepsilon - \bx_0$ \\
Noise schedule & Required & Not needed \\
$t$ distribution & Uniform $\{1,...,T\}$ & Uniform $[0,1]$ \\
Sampling & SDE / DDIM ODE, 50+ steps & Euler ODE, 10–50 steps \\
\bottomrule
\end{tabular}
\end{center}

\section{Reflow: Path Straightening}

Iterative straightening~\cite{liu2022rf}:
\begin{enumerate}
  \item Train $v_\theta^{(0)}$ on random $(x_0, \epsilon)$ pairs
  \item Generate coupled pairs: run ODE to get $\hat{x}_0 = \text{ODE}(\epsilon)$; pair $(\hat{x}_0, \epsilon)$
  \item Train $v_\theta^{(1)}$ on re-paired data — paths are straighter
  \item Repeat — converges to straight paths — 1-step generation becomes feasible
\end{enumerate}

\section{Production: Stable Diffusion 3 and FLUX}

Esser et al.~\cite{esser2024sd3} combine rectified flow with a Multimodal Diffusion Transformer
(MM-DiT): separate transformer streams for text and image tokens that interact via joint attention.
Text conditioning uses CLIP + T5-XXL encoders.

FLUX~\cite{bfl2024flux} extends this with a hybrid DiT architecture and achieves state-of-the-art
open-weight text-to-image quality.

\section*{References}
\begin{thebibliography}{9}
\bibitem{liu2022rf}
  Liu, X., Gong, C., \& Liu, Q. (2022). Flow Straight and Fast: Learning to Generate
  and Transfer Data with Rectified Flow. arXiv:2209.03003.

\bibitem{lipman2022fm}
  Lipman, Y., Chen, R. T. Q., Ben-Hamu, H., Nickel, M., \& Le, M. (2022).
  Flow Matching for Generative Modeling. arXiv:2210.02747.

\bibitem{albergo2022si}
  Albergo, M. S., \& Vanden-Eijnden, E. (2022). Building Normalizing Flows with
  Stochastic Interpolants. arXiv:2209.15571.

\bibitem{esser2024sd3}
  Esser, P., et al. (2024). Scaling Rectified Flow Transformers for High-Resolution
  Image Synthesis (Stable Diffusion 3). arXiv:2403.03206.

\bibitem{bfl2024flux}
  Black Forest Labs (2024). FLUX.1. arXiv:2409.07565.

\bibitem{ma2024sit}
  Ma, N., et al. (2024). SiT: Scalable Interpolant Transformers. arXiv:2401.08740.

\bibitem{song2023consistency}
  Song, Y., et al. (2023). Consistency Models. ICML 2023. arXiv:2303.01469.
\end{thebibliography}

\end{document}
